% !TEX root = ../../../main.tex

\subsection{特征转换与训练数据生成}

在完成 raw dataset 采集后，系统需要将原始样本转换为模型可直接读取的训练数据。
自采样本保存的是连续帧中的手部关键点、姿态关键点和检测状态等原始结构化数据，而模型训练需要的是形状固定、类别明确、数值类型统一的特征样本。
因此，特征转换过程的作用就是将 raw dataset 中按类别保存的样本文件，批量转换为 $30\times166$ 的特征数组，并按照手势类别生成训练标签。
该步骤处于样本采集和模型训练之间，既要保留样本类别语义，又要将不同帧中的关键点整理为统一的数据矩阵。
实时单词识别特征转换流程如图~\ref{fig:chapter06-realtime-feature-convert-flow} 所示。

特征转换脚本首先扫描 raw dataset 根目录。
raw dataset 按手势类别分目录保存，因此脚本会依次遍历每个类别目录，并读取目录下的 sample\_XXX.npz 样本文件。
每读取一个 raw 样本，系统都会检查其帧数是否符合窗口长度要求。
当前实时单词识别模型要求每个样本包含 30 帧，如果某个样本帧数不符合要求，转换流程会跳过该样本并记录失败原因。
这种检查可以避免异常采集片段进入训练集，从而降低后续训练阶段因输入形状不一致导致失败的风险。

单个 raw 样本转换时，系统会逐帧构造 166 维特征。
每一帧先从左右手槽位中读取 hand world landmark 和手部存在状态：若某只手存在，则根据该手的 21 个三维关键点构造 78 维手型特征；
若某只手不存在，则该手对应位置填充为 0。
随后，系统根据姿态关键点构造 10 维姿态辅助特征。
每一帧的左手特征、右手特征和姿态辅助特征按固定顺序拼接，连续 30 帧再按时间顺序堆叠为一个样本。
因此，转换后的每个样本均具有 $30\times166$ 的固定形状，可以直接作为一维时序分类模型的输入。

\WhuAutoFigure{0.82\textwidth}{0.36\textheight}{figures/chapter06/realtime-feature-convert-flow.png}{实时单词识别特征转换流程图}{fig:chapter06-realtime-feature-convert-flow}

表~\ref{tab:chapter06-realtime-feature-conversion-io} 概括了特征转换阶段的输入、处理和输出。
其中，raw 样本仍然保留较完整的检测结果，feature 样本则只保留模型训练所需的数值特征。
这种分层保存方式使系统既可以回溯原始采集数据，也可以直接复用已经转换好的训练数据。

\begin{table}[H]
  \centering
  \caption{实时单词识别特征转换输入输出表}
  \label{tab:chapter06-realtime-feature-conversion-io}
  \begin{tabular}{L{0.18\textwidth} L{0.30\textwidth} L{0.42\textwidth}}
    \toprule
    环节 & 文件或数据 & 说明 \\
    \midrule
    输入目录 & \path{dataset_raw_phone_10fps/<label>/sample_XXX.npz} & 按手势标签分目录保存移动端采集得到的 raw 样本 \\
    帧数检查 & 30 帧窗口 & 只接受满足窗口长度要求的样本，异常样本记录失败原因并跳过 \\
    单帧特征 & 166 维向量 & 由左手 78 维、右手 78 维和姿态辅助 10 维特征拼接得到 \\
    输出目录 & \path{data_processed_arm_pose_10fps/<label>/sample_XXX.npy} & 保存可直接用于模型训练的定长特征样本 \\
    转换摘要 & 扫描数、转换数、跳过数和失败项 & 返回给管理端用于展示转换结果和排查异常样本 \\
    \bottomrule
  \end{tabular}
\end{table}

从工程实现角度看，特征转换并不是简单的文件格式改写。
系统在转换过程中会进行数据可用性检查、缺失关键点填充、手部尺度归一化和姿态相对化处理。
左右手不存在时使用零向量占位，可以保证样本维度稳定；
姿态点不可用或肩部尺度过小时也会采用默认值保护，避免异常数值传播到训练数据中。
完成转换后，管理端可以基于转换摘要判断本次数据集是否已经具备训练条件，并进一步触发固定词表模型训练流程。
